\begin{table}[htbp]
\centering
\caption{χ² Breakdown per Dataset (Detailed)}
\label{tab:χ²_breakdown_per_dataset_(detailed)}
\begin{tabular}{lccccccc}
\hline
Dataset & N_points & Reference & LCDM & χ²/ndf & PBUF & χ²/ndf & Total χ² \\
\hline
Total & 0 & All datasets & 0.00 & 0.000 & 0.00 & 0.000 &  \\
\hline
\end{tabular}
\begin{tablenotes}
\small
\item[1] χ² values represent contribution from each dataset to the total likelihood
\item[2] χ²/ndf shows reduced chi-squared per dataset (χ² divided by number of data points)
\item[3] N_points indicates the number of data points in each dataset
\item[4] References indicate the primary data source for each dataset
\item[5] CMB: Planck 2018 distance priors from acoustic peaks
\item[6] SN: Type Ia supernovae from Pantheon+ compilation
\item[7] BAO: Baryon acoustic oscillation measurements
\item[8] CC: Cosmic chronometers from differential galaxy ages
\item[9] RSD: Redshift-space distortion growth measurements
\end{tablenotes}
\end{table}